% =========================
% CHAPITRE 1 : CONTEXTE GENERAL
% Fichier a inclure depuis le rapport principal via % =========================
% CHAPITRE 1 : CONTEXTE GENERAL
% Fichier a inclure depuis le rapport principal via % =========================
% CHAPITRE 1 : CONTEXTE GENERAL
% Fichier a inclure depuis le rapport principal via % =========================
% CHAPITRE 1 : CONTEXTE GENERAL
% Fichier a inclure depuis le rapport principal via \input{chapters/chapitre1_contexte_general}
% =========================

\chapter{Contexte general}
\label{chap:contexte_general}

\section{Introduction}
Le diabete sucré constitue aujourd'hui l'une des maladies chroniques les plus repandues dans le monde.
Parmi ses complications, la retinopathie diabetique (RD) est particulierement critique car elle peut evoluer
silencieusement pendant plusieurs annees avant l'apparition de symptomes visuels severes.
Sans depistage precoce, cette pathologie peut conduire a une perte de vision irreversible, voire a la cecite.

Sur le plan epidemiologique, la charge de morbidite est considerable. Les donnees de reference
internationales (CDC/OMS, reprises dans les ressources Kaggle de depistage de la RD) indiquent
que plus de 347 millions de personnes vivent avec le diabete dans le monde,
qu'environ 29.1 millions de personnes sont diabetiques aux Etats-Unis,
que la retinopathie diabetique concerne plus de 93 millions de personnes,
et qu'environ 40\% a 45\% des patients diabetiques presentent un certain stade de RD.
Ces chiffres confirment que la RD constitue un enjeu majeur de sante publique,
notamment dans les populations actives.
\textit{Sources : OMS (2024), CDC (2026), NCBI StatPearls (2023).}

Dans les structures de soins, la prise en charge repose encore majoritairement sur une lecture manuelle des images
de fond d'oeil par des specialistes. Cette approche est fiable, mais elle devient difficile a generaliser a grande
échelle en raison du volume de patients, des delais et du manque d'ophtalmologues dans certaines regions.

Dans ce cadre, l'intelligence artificielle (IA) apparait comme une solution d'assistance clinique pertinente :
elle permet d'automatiser une pre-analyse des images, de prioriser les cas urgents et de soutenir la decision
medicale sans remplacer l'expertise du medecin.

\section{Definition de la maladie}
\subsection{Diabete : rappel general}
Le diabete est une maladie metabolique caracterisee par une hyperglycemie chronique.
Cette hyperglycemie peut resulter d'un deficit de secretion d'insuline, d'une resistance a l'insuline,
ou d'une combinaison des deux. A long terme, l'exces de glucose sanguin provoque des atteintes vasculaires
et nerveuses qui touchent plusieurs organes (oeil, rein, coeur, systeme nerveux).

\subsection{Types de diabete}
\begin{itemize}
    \item \textbf{Diabete de type 1} : maladie auto-immune, destruction des cellules beta pancreatiques,
    dependance a l'insuline.
    \item \textbf{Diabete de type 2} : forme la plus frequente, associee a l'insulino-resistance et au mode de vie.
    \item \textbf{Diabete gestationnel} : hyperglycemie detectee pendant la grossesse.
    \item \textbf{Autres formes} : causes genetiques, endocriniennes ou medicamenteuses.
\end{itemize}

\subsection{Anatomie de l'oeil et zones d'interet}
Pour comprendre la retinopathie diabetique, il est essentiel de rappeler les structures oculaires impliquees
dans le diagnostic :
\begin{itemize}
    \item \textbf{Retine} : tissu neurosensoriel convertissant la lumiere en influx nerveux.
    \item \textbf{Macula} : zone centrale de la vision fine, critique pour la lecture et la precision visuelle.
    \item \textbf{Fovea} : centre de la macula, point de plus haute acuite visuelle.
    \item \textbf{Papille optique} : zone d'emergence du nerf optique.
    \item \textbf{Vaisseaux retiniens} : structures les plus affectees par les complications diabetiques.
\end{itemize}

\begin{figure}[H]
    \centering
    \includegraphics[width=0.72\textwidth]{images/ch1/anatomie_oeil.png}
    \caption{Anatomie simplifiee de l'oeil et principales zones d'interet pour le depistage de la RD.}
    \label{fig:anatomie_oeil_ch1}
\end{figure}

\subsection{Retinopathie diabetique : mecanisme et progression}
La retinopathie diabetique (RD) est une microangiopathie chronique de la retine associee
a un diabete mal controle ou ancien. L'hyperglycemie prolongee, associee a d'autres facteurs
de risque (hypertension arterielle, dyslipidemie, surcharge ponderale), fragilise l'endothelium
vasculaire et provoque des anomalies progressives de la microcirculation retinienne.

\subsubsection*{Lesions retiniennes elementaires}
Dans les images du fond d'oeil, la RD est principalement reconnue par la presence de lesions
vasculaires caracteristiques :
\begin{itemize}
    \item microanevrismes ;
    \item hemorragies intraretiniennes ;
    \item exsudats lipidiques ;
    \item exsudats cotonneux (zones d'ischemie nerveuse) ;
    \item anomalies veineuses (beading, irregularites) et ischemie ;
    \item neovascularisation dans les stades avances.
\end{itemize}

\subsubsection*{Stades cliniques : forme non proliferante et proliferante}
Sur le plan clinique, la progression suit deux grandes phases.

\paragraph{1) Retinopathie diabetique non proliferante (NPDR)}
Il s'agit du stade initial, le plus frequent au debut de l'evolution. Les petits vaisseaux
retiniens deviennent permeables, avec fuite de sang et de liquide. Cette fuite peut entrainer
un oedeme maculaire, cause majeure de baisse visuelle chez les personnes diabetiques.
Dans les stades non proliferants, la vision peut rester longtemps peu alteree, ce qui retarde
la consultation et renforce l'importance du depistage periodique.

\paragraph{2) Retinopathie diabetique proliferante (PDR)}
Lorsque l'ischemie retinienne devient importante (vaisseaux occlus), la retine declenche
une neovascularisation pathologique. Ces nouveaux vaisseaux sont fragiles, inefficaces
et susceptibles de saigner dans le vitre, provoquant flou visuel et corps flottants.
Le stade proliferant est associe a un risque eleve de complications severes.

\subsubsection*{Complications majeures des formes avancees}
\begin{itemize}
    \item \textbf{Hemorragie du vitre} : baisse visuelle brutale, parfois severe.
    \item \textbf{Fibrose retinienne} : formation de tissu cicatriciel pouvant deformer la retine.
    \item \textbf{Decollement tractionnel de retine} : urgence fonctionnelle menacant la vision centrale.
    \item \textbf{Glaucome neovasculaire} : obstruction de l'angle irido-corneen et elevation rapide
    de la pression intraoculaire.
\end{itemize}

\subsubsection*{Classification simplifiee exploitee dans le projet}
Pour la modelisation IA et l'evaluation, on retient une gradation en cinq niveaux :
\begin{enumerate}
    \item \textbf{Absence de RD} ;
    \item \textbf{RD non proliferante legere} ;
    \item \textbf{RD non proliferante moderee} ;
    \item \textbf{RD non proliferante severe} ;
    \item \textbf{RD proliferante} (risque eleve de perte visuelle).
\end{enumerate}

\begin{figure}[H]
    \centering
    \includegraphics[width=0.85\textwidth]{images/ch1/stades_retinopathie.png}
    \caption{Exemple visuel des stades de progression de la retinopathie diabetique.}
    \label{fig:stades_rd_ch1}
\end{figure}

\paragraph{Lien avec le depistage automatise}
Les informations precedentes montrent qu'une detection rapide des lesions est essentielle
pour eviter la progression vers les formes proliferantes. C'est exactement l'objectif
des systemes d'aide au depistage automatise : analyser les images de fond d'oeil,
estimer un grade de severite et orienter en priorite les patients a risque eleve.

\section{Methodes de diagnostic existantes}
\subsection{Examen clinique du fond d'oeil}
L'examen du fond d'oeil est la methode standard pour observer les lesions retiniennes.
Il peut etre realise avec ou sans dilatation pupillaire, selon le contexte clinique.

\subsection{Retinographie couleur}
La retinographie permet de capturer des images numeriques du fond d'oeil.
Elle facilite l'archivage, la tele-expertise et le suivi longitudinal des patients.

\subsection{Tomographie en coherence optique (OCT)}
L'OCT fournit des coupes de la retine a haute resolution.
Cette technique est particulierement utile pour detecter l'oedeme maculaire diabetique.

\subsection{Angiographie retinienne}
L'angiographie a la fluoresceine visualise la perfusion vasculaire et les zones de fuite.
Bien qu'informative, elle est plus invasive et plus couteuse.

\subsection{Tele-ophtalmologie}
La tele-ophtalmologie permet de capturer localement les images puis de les transmettre
a un expert a distance. Cette approche augmente l'accessibilite au depistage.

\begin{table}[H]
\centering
\caption{Comparaison des principales methodes de diagnostic de la retinopathie diabetique.}
\label{tab:comparaison_diagnostic_ch1}
\begin{tabular}{|p{3.8cm}|p{3.2cm}|p{3.6cm}|p{3.6cm}|}
\hline
\textbf{Methode} & \textbf{Atouts} & \textbf{Limites} & \textbf{Contexte d'usage} \\
\hline
Fond d'oeil clinique & Methode de reference, pratique courante & Dependance a l'expertise humaine & Depistage et suivi \\
\hline
Retinographie & Trace numerique, partage a distance & Qualite variable selon acquisition & Campagnes de depistage \\
\hline
OCT & Tres haute precision structurelle & Cout eleve, equipement specialise & Oedeme maculaire, cas complexes \\
\hline
Angiographie & Analyse vasculaire detaillee & Invasive, moins accessible & Exploration avancee \\
\hline
Tele-ophtalmologie & Couvre les zones eloignees & Delai d'interpretation possible & Reseaux de soins distribues \\
\hline
\end{tabular}
\end{table}

\section{Prediction du diabete par l'intelligence artificielle}
Dans le contexte actuel, le depistage manuel de la retinopathie diabetique reste efficace
mais il est limite par la charge de travail des specialistes, la variabilite de lecture,
les delais de restitution et les inegalites d'acces aux soins.
L'intelligence artificielle repond a cette problematique en fournissant une aide rapide
et standardisee pour analyser les images du fond d'oeil, estimer un grade de severite,
et prioriser les cas urgents.

L'objectif n'est pas de remplacer l'ophtalmologue, mais de l'assister dans la decision,
de reduire le retard diagnostique et d'ameliorer la qualite globale du depistage.

\subsubsection*{1.4.7 Solution avec IA et concept de l'application}
La solution proposee dans ce projet repose sur une application d'aide au depistage
qui combine une interface clinique (centre/medecin), un backend de gestion des examens
et un service IA d'analyse d'images retiniennes.

Le concept de l'application est le suivant :
\begin{enumerate}
    \item acquisition de l'image du fond d'oeil par le centre ;
    \item envoi securise de l'image vers le serveur d'analyse ;
    \item prediction automatique du grade de retinopathie ;
    \item generation d'une explication visuelle (Grad-CAM) ;
    \item notification du medecin, avec priorisation des cas severes.
\end{enumerate}

Cette architecture permet de reduire les delais de depistage,
de standardiser la pre-analyse et d'ameliorer la coordination entre les acteurs de soin.

\section{Methodologie du projet (Scrum)}
Pour cadrer le developpement de l'application, nous avons adopte une methode agile de type Scrum,
adaptee aux projets iteratifs en intelligence artificielle et en developpement logiciel.

Le travail a ete organise en sprints courts, avec des objectifs incrementaux :
\begin{itemize}
    \item preparation et exploration des donnees ;
    \item experimentation et amelioration des modeles ;
    \item integration backend/frontend/service IA ;
    \item validation fonctionnelle et tests de bout en bout.
\end{itemize}

Scrum a apporte trois benefices principaux dans ce projet :
\begin{itemize}
    \item \textbf{adaptabilite} : ajustement rapide des priorites selon les resultats experimentaux ;
    \item \textbf{visibilite} : suivi continu des taches et des livrables ;
    \item \textbf{qualite} : integration progressive et verification reguliere des fonctionnalites.
\end{itemize}

\section{Partie statistique}
\subsection{a) Dans le monde}
Les indicateurs epidemiologiques internationaux confirment l'ampleur du probleme :
\begin{itemize}
    \item plus de 347 millions de personnes vivent avec le diabete dans le monde ;
    \item la retinopathie diabetique concernerait plus de 93 millions de personnes ;
    \item environ 40\% a 45\% des personnes diabetiques presentent un stade de retinopathie.
\end{itemize}

Ces chiffres montrent que le depistage de la RD n'est pas uniquement un enjeu clinique,
mais aussi un enjeu de sante publique globale, en particulier dans les regions ou
l'acces aux specialistes reste limite.
\textit{Sources : WHO Fact Sheet Diabetes (2024), IDF Diabetes Atlas (2025).}

\subsection{b) En Tunisie}
En Tunisie, le diabete represente egalement un probleme majeur de sante publique,
avec une prevalence estimee autour de 16\% chez les adultes (20--79 ans) selon la serie
IDF/World Bank (2024). Cette progression s'accompagne d'un risque accru de complications
chroniques, dont la retinopathie diabetique.

Dans ce contexte, le depistage precoce reste un defi operationnel, notamment en raison :
\begin{itemize}
    \item de la concentration des specialistes dans les grands centres urbains ;
    \item des delais de consultation en ophtalmologie ;
    \item de l'augmentation continue du nombre de patients diabetiques.
\end{itemize}

L'integration d'une solution d'aide au depistage basee sur l'IA peut ainsi contribuer
a reduire les delais d'orientation et a ameliorer l'acces au diagnostic sur le territoire.
\textit{Source : World Bank Data -- Indicator SH.STA.DIAB.ZS (Tunisia, 2024).}

\section*{References web utilisees}
\addcontentsline{toc}{section}{References web utilisees}
\begin{itemize}
    \item World Health Organization (WHO). \textit{Diabetes -- Fact sheet}.\newline
    \url{https://www.who.int/news-room/fact-sheets/detail/diabetes}
    \item Centers for Disease Control and Prevention (CDC). \textit{Diabetes Basics}.\newline
    \url{https://www.cdc.gov/diabetes/about/index.html}
    \item International Diabetes Federation (IDF). \textit{IDF Diabetes Atlas}.\newline
    \url{https://diabetesatlas.org/}
    \item World Bank Data. \textit{Diabetes prevalence (\% of population ages 20 to 79) -- Tunisia}.\newline
    \url{https://data.worldbank.org/indicator/SH.STA.DIAB.ZS?locations=TN}
    \item NCBI Bookshelf (StatPearls). \textit{Diabetic Retinopathy}.\newline
    \url{https://www.ncbi.nlm.nih.gov/books/NBK560805/}
\end{itemize}

\section{Conclusion}
Ce chapitre a presente le contexte general du projet, en partant des bases medicales du diabete
et de la retinopathie diabetique, puis en decrivant les methodes diagnostiques actuelles et leurs limites.
L'analyse des contraintes terrain justifie l'interet d'une solution numerique intelligente.

La suite du rapport s'orientera vers l'etat de l'art des approches IA en imagerie retinienne,
les choix methodologiques retenus et la conception technique de la solution proposee.

% =========================

\chapter{Contexte general}
\label{chap:contexte_general}

\section{Introduction}
Le diabete sucré constitue aujourd'hui l'une des maladies chroniques les plus repandues dans le monde.
Parmi ses complications, la retinopathie diabetique (RD) est particulierement critique car elle peut evoluer
silencieusement pendant plusieurs annees avant l'apparition de symptomes visuels severes.
Sans depistage precoce, cette pathologie peut conduire a une perte de vision irreversible, voire a la cecite.

Sur le plan epidemiologique, la charge de morbidite est considerable. Les donnees de reference
internationales (CDC/OMS, reprises dans les ressources Kaggle de depistage de la RD) indiquent
que plus de 347 millions de personnes vivent avec le diabete dans le monde,
qu'environ 29.1 millions de personnes sont diabetiques aux Etats-Unis,
que la retinopathie diabetique concerne plus de 93 millions de personnes,
et qu'environ 40\% a 45\% des patients diabetiques presentent un certain stade de RD.
Ces chiffres confirment que la RD constitue un enjeu majeur de sante publique,
notamment dans les populations actives.
\textit{Sources : OMS (2024), CDC (2026), NCBI StatPearls (2023).}

Dans les structures de soins, la prise en charge repose encore majoritairement sur une lecture manuelle des images
de fond d'oeil par des specialistes. Cette approche est fiable, mais elle devient difficile a generaliser a grande
échelle en raison du volume de patients, des delais et du manque d'ophtalmologues dans certaines regions.

Dans ce cadre, l'intelligence artificielle (IA) apparait comme une solution d'assistance clinique pertinente :
elle permet d'automatiser une pre-analyse des images, de prioriser les cas urgents et de soutenir la decision
medicale sans remplacer l'expertise du medecin.

\section{Definition de la maladie}
\subsection{Diabete : rappel general}
Le diabete est une maladie metabolique caracterisee par une hyperglycemie chronique.
Cette hyperglycemie peut resulter d'un deficit de secretion d'insuline, d'une resistance a l'insuline,
ou d'une combinaison des deux. A long terme, l'exces de glucose sanguin provoque des atteintes vasculaires
et nerveuses qui touchent plusieurs organes (oeil, rein, coeur, systeme nerveux).

\subsection{Types de diabete}
\begin{itemize}
    \item \textbf{Diabete de type 1} : maladie auto-immune, destruction des cellules beta pancreatiques,
    dependance a l'insuline.
    \item \textbf{Diabete de type 2} : forme la plus frequente, associee a l'insulino-resistance et au mode de vie.
    \item \textbf{Diabete gestationnel} : hyperglycemie detectee pendant la grossesse.
    \item \textbf{Autres formes} : causes genetiques, endocriniennes ou medicamenteuses.
\end{itemize}

\subsection{Anatomie de l'oeil et zones d'interet}
Pour comprendre la retinopathie diabetique, il est essentiel de rappeler les structures oculaires impliquees
dans le diagnostic :
\begin{itemize}
    \item \textbf{Retine} : tissu neurosensoriel convertissant la lumiere en influx nerveux.
    \item \textbf{Macula} : zone centrale de la vision fine, critique pour la lecture et la precision visuelle.
    \item \textbf{Fovea} : centre de la macula, point de plus haute acuite visuelle.
    \item \textbf{Papille optique} : zone d'emergence du nerf optique.
    \item \textbf{Vaisseaux retiniens} : structures les plus affectees par les complications diabetiques.
\end{itemize}

\begin{figure}[H]
    \centering
    \includegraphics[width=0.72\textwidth]{images/ch1/anatomie_oeil.png}
    \caption{Anatomie simplifiee de l'oeil et principales zones d'interet pour le depistage de la RD.}
    \label{fig:anatomie_oeil_ch1}
\end{figure}

\subsection{Retinopathie diabetique : mecanisme et progression}
La retinopathie diabetique (RD) est une microangiopathie chronique de la retine associee
a un diabete mal controle ou ancien. L'hyperglycemie prolongee, associee a d'autres facteurs
de risque (hypertension arterielle, dyslipidemie, surcharge ponderale), fragilise l'endothelium
vasculaire et provoque des anomalies progressives de la microcirculation retinienne.

\subsubsection*{Lesions retiniennes elementaires}
Dans les images du fond d'oeil, la RD est principalement reconnue par la presence de lesions
vasculaires caracteristiques :
\begin{itemize}
    \item microanevrismes ;
    \item hemorragies intraretiniennes ;
    \item exsudats lipidiques ;
    \item exsudats cotonneux (zones d'ischemie nerveuse) ;
    \item anomalies veineuses (beading, irregularites) et ischemie ;
    \item neovascularisation dans les stades avances.
\end{itemize}

\subsubsection*{Stades cliniques : forme non proliferante et proliferante}
Sur le plan clinique, la progression suit deux grandes phases.

\paragraph{1) Retinopathie diabetique non proliferante (NPDR)}
Il s'agit du stade initial, le plus frequent au debut de l'evolution. Les petits vaisseaux
retiniens deviennent permeables, avec fuite de sang et de liquide. Cette fuite peut entrainer
un oedeme maculaire, cause majeure de baisse visuelle chez les personnes diabetiques.
Dans les stades non proliferants, la vision peut rester longtemps peu alteree, ce qui retarde
la consultation et renforce l'importance du depistage periodique.

\paragraph{2) Retinopathie diabetique proliferante (PDR)}
Lorsque l'ischemie retinienne devient importante (vaisseaux occlus), la retine declenche
une neovascularisation pathologique. Ces nouveaux vaisseaux sont fragiles, inefficaces
et susceptibles de saigner dans le vitre, provoquant flou visuel et corps flottants.
Le stade proliferant est associe a un risque eleve de complications severes.

\subsubsection*{Complications majeures des formes avancees}
\begin{itemize}
    \item \textbf{Hemorragie du vitre} : baisse visuelle brutale, parfois severe.
    \item \textbf{Fibrose retinienne} : formation de tissu cicatriciel pouvant deformer la retine.
    \item \textbf{Decollement tractionnel de retine} : urgence fonctionnelle menacant la vision centrale.
    \item \textbf{Glaucome neovasculaire} : obstruction de l'angle irido-corneen et elevation rapide
    de la pression intraoculaire.
\end{itemize}

\subsubsection*{Classification simplifiee exploitee dans le projet}
Pour la modelisation IA et l'evaluation, on retient une gradation en cinq niveaux :
\begin{enumerate}
    \item \textbf{Absence de RD} ;
    \item \textbf{RD non proliferante legere} ;
    \item \textbf{RD non proliferante moderee} ;
    \item \textbf{RD non proliferante severe} ;
    \item \textbf{RD proliferante} (risque eleve de perte visuelle).
\end{enumerate}

\begin{figure}[H]
    \centering
    \includegraphics[width=0.85\textwidth]{images/ch1/stades_retinopathie.png}
    \caption{Exemple visuel des stades de progression de la retinopathie diabetique.}
    \label{fig:stades_rd_ch1}
\end{figure}

\paragraph{Lien avec le depistage automatise}
Les informations precedentes montrent qu'une detection rapide des lesions est essentielle
pour eviter la progression vers les formes proliferantes. C'est exactement l'objectif
des systemes d'aide au depistage automatise : analyser les images de fond d'oeil,
estimer un grade de severite et orienter en priorite les patients a risque eleve.

\section{Methodes de diagnostic existantes}
\subsection{Examen clinique du fond d'oeil}
L'examen du fond d'oeil est la methode standard pour observer les lesions retiniennes.
Il peut etre realise avec ou sans dilatation pupillaire, selon le contexte clinique.

\subsection{Retinographie couleur}
La retinographie permet de capturer des images numeriques du fond d'oeil.
Elle facilite l'archivage, la tele-expertise et le suivi longitudinal des patients.

\subsection{Tomographie en coherence optique (OCT)}
L'OCT fournit des coupes de la retine a haute resolution.
Cette technique est particulierement utile pour detecter l'oedeme maculaire diabetique.

\subsection{Angiographie retinienne}
L'angiographie a la fluoresceine visualise la perfusion vasculaire et les zones de fuite.
Bien qu'informative, elle est plus invasive et plus couteuse.

\subsection{Tele-ophtalmologie}
La tele-ophtalmologie permet de capturer localement les images puis de les transmettre
a un expert a distance. Cette approche augmente l'accessibilite au depistage.

\begin{table}[H]
\centering
\caption{Comparaison des principales methodes de diagnostic de la retinopathie diabetique.}
\label{tab:comparaison_diagnostic_ch1}
\begin{tabular}{|p{3.8cm}|p{3.2cm}|p{3.6cm}|p{3.6cm}|}
\hline
\textbf{Methode} & \textbf{Atouts} & \textbf{Limites} & \textbf{Contexte d'usage} \\
\hline
Fond d'oeil clinique & Methode de reference, pratique courante & Dependance a l'expertise humaine & Depistage et suivi \\
\hline
Retinographie & Trace numerique, partage a distance & Qualite variable selon acquisition & Campagnes de depistage \\
\hline
OCT & Tres haute precision structurelle & Cout eleve, equipement specialise & Oedeme maculaire, cas complexes \\
\hline
Angiographie & Analyse vasculaire detaillee & Invasive, moins accessible & Exploration avancee \\
\hline
Tele-ophtalmologie & Couvre les zones eloignees & Delai d'interpretation possible & Reseaux de soins distribues \\
\hline
\end{tabular}
\end{table}

\section{Prediction du diabete par l'intelligence artificielle}
Dans le contexte actuel, le depistage manuel de la retinopathie diabetique reste efficace
mais il est limite par la charge de travail des specialistes, la variabilite de lecture,
les delais de restitution et les inegalites d'acces aux soins.
L'intelligence artificielle repond a cette problematique en fournissant une aide rapide
et standardisee pour analyser les images du fond d'oeil, estimer un grade de severite,
et prioriser les cas urgents.

L'objectif n'est pas de remplacer l'ophtalmologue, mais de l'assister dans la decision,
de reduire le retard diagnostique et d'ameliorer la qualite globale du depistage.

\subsubsection*{1.4.7 Solution avec IA et concept de l'application}
La solution proposee dans ce projet repose sur une application d'aide au depistage
qui combine une interface clinique (centre/medecin), un backend de gestion des examens
et un service IA d'analyse d'images retiniennes.

Le concept de l'application est le suivant :
\begin{enumerate}
    \item acquisition de l'image du fond d'oeil par le centre ;
    \item envoi securise de l'image vers le serveur d'analyse ;
    \item prediction automatique du grade de retinopathie ;
    \item generation d'une explication visuelle (Grad-CAM) ;
    \item notification du medecin, avec priorisation des cas severes.
\end{enumerate}

Cette architecture permet de reduire les delais de depistage,
de standardiser la pre-analyse et d'ameliorer la coordination entre les acteurs de soin.

\section{Methodologie du projet (Scrum)}
Pour cadrer le developpement de l'application, nous avons adopte une methode agile de type Scrum,
adaptee aux projets iteratifs en intelligence artificielle et en developpement logiciel.

Le travail a ete organise en sprints courts, avec des objectifs incrementaux :
\begin{itemize}
    \item preparation et exploration des donnees ;
    \item experimentation et amelioration des modeles ;
    \item integration backend/frontend/service IA ;
    \item validation fonctionnelle et tests de bout en bout.
\end{itemize}

Scrum a apporte trois benefices principaux dans ce projet :
\begin{itemize}
    \item \textbf{adaptabilite} : ajustement rapide des priorites selon les resultats experimentaux ;
    \item \textbf{visibilite} : suivi continu des taches et des livrables ;
    \item \textbf{qualite} : integration progressive et verification reguliere des fonctionnalites.
\end{itemize}

\section{Partie statistique}
\subsection{a) Dans le monde}
Les indicateurs epidemiologiques internationaux confirment l'ampleur du probleme :
\begin{itemize}
    \item plus de 347 millions de personnes vivent avec le diabete dans le monde ;
    \item la retinopathie diabetique concernerait plus de 93 millions de personnes ;
    \item environ 40\% a 45\% des personnes diabetiques presentent un stade de retinopathie.
\end{itemize}

Ces chiffres montrent que le depistage de la RD n'est pas uniquement un enjeu clinique,
mais aussi un enjeu de sante publique globale, en particulier dans les regions ou
l'acces aux specialistes reste limite.
\textit{Sources : WHO Fact Sheet Diabetes (2024), IDF Diabetes Atlas (2025).}

\subsection{b) En Tunisie}
En Tunisie, le diabete represente egalement un probleme majeur de sante publique,
avec une prevalence estimee autour de 16\% chez les adultes (20--79 ans) selon la serie
IDF/World Bank (2024). Cette progression s'accompagne d'un risque accru de complications
chroniques, dont la retinopathie diabetique.

Dans ce contexte, le depistage precoce reste un defi operationnel, notamment en raison :
\begin{itemize}
    \item de la concentration des specialistes dans les grands centres urbains ;
    \item des delais de consultation en ophtalmologie ;
    \item de l'augmentation continue du nombre de patients diabetiques.
\end{itemize}

L'integration d'une solution d'aide au depistage basee sur l'IA peut ainsi contribuer
a reduire les delais d'orientation et a ameliorer l'acces au diagnostic sur le territoire.
\textit{Source : World Bank Data -- Indicator SH.STA.DIAB.ZS (Tunisia, 2024).}

\section*{References web utilisees}
\addcontentsline{toc}{section}{References web utilisees}
\begin{itemize}
    \item World Health Organization (WHO). \textit{Diabetes -- Fact sheet}.\newline
    \url{https://www.who.int/news-room/fact-sheets/detail/diabetes}
    \item Centers for Disease Control and Prevention (CDC). \textit{Diabetes Basics}.\newline
    \url{https://www.cdc.gov/diabetes/about/index.html}
    \item International Diabetes Federation (IDF). \textit{IDF Diabetes Atlas}.\newline
    \url{https://diabetesatlas.org/}
    \item World Bank Data. \textit{Diabetes prevalence (\% of population ages 20 to 79) -- Tunisia}.\newline
    \url{https://data.worldbank.org/indicator/SH.STA.DIAB.ZS?locations=TN}
    \item NCBI Bookshelf (StatPearls). \textit{Diabetic Retinopathy}.\newline
    \url{https://www.ncbi.nlm.nih.gov/books/NBK560805/}
\end{itemize}

\section{Conclusion}
Ce chapitre a presente le contexte general du projet, en partant des bases medicales du diabete
et de la retinopathie diabetique, puis en decrivant les methodes diagnostiques actuelles et leurs limites.
L'analyse des contraintes terrain justifie l'interet d'une solution numerique intelligente.

La suite du rapport s'orientera vers l'etat de l'art des approches IA en imagerie retinienne,
les choix methodologiques retenus et la conception technique de la solution proposee.

% =========================

\chapter{Contexte general}
\label{chap:contexte_general}

\section{Introduction}
Le diabete sucré constitue aujourd'hui l'une des maladies chroniques les plus repandues dans le monde.
Parmi ses complications, la retinopathie diabetique (RD) est particulierement critique car elle peut evoluer
silencieusement pendant plusieurs annees avant l'apparition de symptomes visuels severes.
Sans depistage precoce, cette pathologie peut conduire a une perte de vision irreversible, voire a la cecite.

Sur le plan epidemiologique, la charge de morbidite est considerable. Les donnees de reference
internationales (CDC/OMS, reprises dans les ressources Kaggle de depistage de la RD) indiquent
que plus de 347 millions de personnes vivent avec le diabete dans le monde,
qu'environ 29.1 millions de personnes sont diabetiques aux Etats-Unis,
que la retinopathie diabetique concerne plus de 93 millions de personnes,
et qu'environ 40\% a 45\% des patients diabetiques presentent un certain stade de RD.
Ces chiffres confirment que la RD constitue un enjeu majeur de sante publique,
notamment dans les populations actives.
\textit{Sources : OMS (2024), CDC (2026), NCBI StatPearls (2023).}

Dans les structures de soins, la prise en charge repose encore majoritairement sur une lecture manuelle des images
de fond d'oeil par des specialistes. Cette approche est fiable, mais elle devient difficile a generaliser a grande
échelle en raison du volume de patients, des delais et du manque d'ophtalmologues dans certaines regions.

Dans ce cadre, l'intelligence artificielle (IA) apparait comme une solution d'assistance clinique pertinente :
elle permet d'automatiser une pre-analyse des images, de prioriser les cas urgents et de soutenir la decision
medicale sans remplacer l'expertise du medecin.

\section{Definition de la maladie}
\subsection{Diabete : rappel general}
Le diabete est une maladie metabolique caracterisee par une hyperglycemie chronique.
Cette hyperglycemie peut resulter d'un deficit de secretion d'insuline, d'une resistance a l'insuline,
ou d'une combinaison des deux. A long terme, l'exces de glucose sanguin provoque des atteintes vasculaires
et nerveuses qui touchent plusieurs organes (oeil, rein, coeur, systeme nerveux).

\subsection{Types de diabete}
\begin{itemize}
    \item \textbf{Diabete de type 1} : maladie auto-immune, destruction des cellules beta pancreatiques,
    dependance a l'insuline.
    \item \textbf{Diabete de type 2} : forme la plus frequente, associee a l'insulino-resistance et au mode de vie.
    \item \textbf{Diabete gestationnel} : hyperglycemie detectee pendant la grossesse.
    \item \textbf{Autres formes} : causes genetiques, endocriniennes ou medicamenteuses.
\end{itemize}

\subsection{Anatomie de l'oeil et zones d'interet}
Pour comprendre la retinopathie diabetique, il est essentiel de rappeler les structures oculaires impliquees
dans le diagnostic :
\begin{itemize}
    \item \textbf{Retine} : tissu neurosensoriel convertissant la lumiere en influx nerveux.
    \item \textbf{Macula} : zone centrale de la vision fine, critique pour la lecture et la precision visuelle.
    \item \textbf{Fovea} : centre de la macula, point de plus haute acuite visuelle.
    \item \textbf{Papille optique} : zone d'emergence du nerf optique.
    \item \textbf{Vaisseaux retiniens} : structures les plus affectees par les complications diabetiques.
\end{itemize}

\begin{figure}[H]
    \centering
    \includegraphics[width=0.72\textwidth]{images/ch1/anatomie_oeil.png}
    \caption{Anatomie simplifiee de l'oeil et principales zones d'interet pour le depistage de la RD.}
    \label{fig:anatomie_oeil_ch1}
\end{figure}

\subsection{Retinopathie diabetique : mecanisme et progression}
La retinopathie diabetique (RD) est une microangiopathie chronique de la retine associee
a un diabete mal controle ou ancien. L'hyperglycemie prolongee, associee a d'autres facteurs
de risque (hypertension arterielle, dyslipidemie, surcharge ponderale), fragilise l'endothelium
vasculaire et provoque des anomalies progressives de la microcirculation retinienne.

\subsubsection*{Lesions retiniennes elementaires}
Dans les images du fond d'oeil, la RD est principalement reconnue par la presence de lesions
vasculaires caracteristiques :
\begin{itemize}
    \item microanevrismes ;
    \item hemorragies intraretiniennes ;
    \item exsudats lipidiques ;
    \item exsudats cotonneux (zones d'ischemie nerveuse) ;
    \item anomalies veineuses (beading, irregularites) et ischemie ;
    \item neovascularisation dans les stades avances.
\end{itemize}

\subsubsection*{Stades cliniques : forme non proliferante et proliferante}
Sur le plan clinique, la progression suit deux grandes phases.

\paragraph{1) Retinopathie diabetique non proliferante (NPDR)}
Il s'agit du stade initial, le plus frequent au debut de l'evolution. Les petits vaisseaux
retiniens deviennent permeables, avec fuite de sang et de liquide. Cette fuite peut entrainer
un oedeme maculaire, cause majeure de baisse visuelle chez les personnes diabetiques.
Dans les stades non proliferants, la vision peut rester longtemps peu alteree, ce qui retarde
la consultation et renforce l'importance du depistage periodique.

\paragraph{2) Retinopathie diabetique proliferante (PDR)}
Lorsque l'ischemie retinienne devient importante (vaisseaux occlus), la retine declenche
une neovascularisation pathologique. Ces nouveaux vaisseaux sont fragiles, inefficaces
et susceptibles de saigner dans le vitre, provoquant flou visuel et corps flottants.
Le stade proliferant est associe a un risque eleve de complications severes.

\subsubsection*{Complications majeures des formes avancees}
\begin{itemize}
    \item \textbf{Hemorragie du vitre} : baisse visuelle brutale, parfois severe.
    \item \textbf{Fibrose retinienne} : formation de tissu cicatriciel pouvant deformer la retine.
    \item \textbf{Decollement tractionnel de retine} : urgence fonctionnelle menacant la vision centrale.
    \item \textbf{Glaucome neovasculaire} : obstruction de l'angle irido-corneen et elevation rapide
    de la pression intraoculaire.
\end{itemize}

\subsubsection*{Classification simplifiee exploitee dans le projet}
Pour la modelisation IA et l'evaluation, on retient une gradation en cinq niveaux :
\begin{enumerate}
    \item \textbf{Absence de RD} ;
    \item \textbf{RD non proliferante legere} ;
    \item \textbf{RD non proliferante moderee} ;
    \item \textbf{RD non proliferante severe} ;
    \item \textbf{RD proliferante} (risque eleve de perte visuelle).
\end{enumerate}

\begin{figure}[H]
    \centering
    \includegraphics[width=0.85\textwidth]{images/ch1/stades_retinopathie.png}
    \caption{Exemple visuel des stades de progression de la retinopathie diabetique.}
    \label{fig:stades_rd_ch1}
\end{figure}

\paragraph{Lien avec le depistage automatise}
Les informations precedentes montrent qu'une detection rapide des lesions est essentielle
pour eviter la progression vers les formes proliferantes. C'est exactement l'objectif
des systemes d'aide au depistage automatise : analyser les images de fond d'oeil,
estimer un grade de severite et orienter en priorite les patients a risque eleve.

\section{Methodes de diagnostic existantes}
\subsection{Examen clinique du fond d'oeil}
L'examen du fond d'oeil est la methode standard pour observer les lesions retiniennes.
Il peut etre realise avec ou sans dilatation pupillaire, selon le contexte clinique.

\subsection{Retinographie couleur}
La retinographie permet de capturer des images numeriques du fond d'oeil.
Elle facilite l'archivage, la tele-expertise et le suivi longitudinal des patients.

\subsection{Tomographie en coherence optique (OCT)}
L'OCT fournit des coupes de la retine a haute resolution.
Cette technique est particulierement utile pour detecter l'oedeme maculaire diabetique.

\subsection{Angiographie retinienne}
L'angiographie a la fluoresceine visualise la perfusion vasculaire et les zones de fuite.
Bien qu'informative, elle est plus invasive et plus couteuse.

\subsection{Tele-ophtalmologie}
La tele-ophtalmologie permet de capturer localement les images puis de les transmettre
a un expert a distance. Cette approche augmente l'accessibilite au depistage.

\begin{table}[H]
\centering
\caption{Comparaison des principales methodes de diagnostic de la retinopathie diabetique.}
\label{tab:comparaison_diagnostic_ch1}
\begin{tabular}{|p{3.8cm}|p{3.2cm}|p{3.6cm}|p{3.6cm}|}
\hline
\textbf{Methode} & \textbf{Atouts} & \textbf{Limites} & \textbf{Contexte d'usage} \\
\hline
Fond d'oeil clinique & Methode de reference, pratique courante & Dependance a l'expertise humaine & Depistage et suivi \\
\hline
Retinographie & Trace numerique, partage a distance & Qualite variable selon acquisition & Campagnes de depistage \\
\hline
OCT & Tres haute precision structurelle & Cout eleve, equipement specialise & Oedeme maculaire, cas complexes \\
\hline
Angiographie & Analyse vasculaire detaillee & Invasive, moins accessible & Exploration avancee \\
\hline
Tele-ophtalmologie & Couvre les zones eloignees & Delai d'interpretation possible & Reseaux de soins distribues \\
\hline
\end{tabular}
\end{table}

\section{Prediction du diabete par l'intelligence artificielle}
Dans le contexte actuel, le depistage manuel de la retinopathie diabetique reste efficace
mais il est limite par la charge de travail des specialistes, la variabilite de lecture,
les delais de restitution et les inegalites d'acces aux soins.
L'intelligence artificielle repond a cette problematique en fournissant une aide rapide
et standardisee pour analyser les images du fond d'oeil, estimer un grade de severite,
et prioriser les cas urgents.

L'objectif n'est pas de remplacer l'ophtalmologue, mais de l'assister dans la decision,
de reduire le retard diagnostique et d'ameliorer la qualite globale du depistage.

\subsubsection*{1.4.7 Solution avec IA et concept de l'application}
La solution proposee dans ce projet repose sur une application d'aide au depistage
qui combine une interface clinique (centre/medecin), un backend de gestion des examens
et un service IA d'analyse d'images retiniennes.

Le concept de l'application est le suivant :
\begin{enumerate}
    \item acquisition de l'image du fond d'oeil par le centre ;
    \item envoi securise de l'image vers le serveur d'analyse ;
    \item prediction automatique du grade de retinopathie ;
    \item generation d'une explication visuelle (Grad-CAM) ;
    \item notification du medecin, avec priorisation des cas severes.
\end{enumerate}

Cette architecture permet de reduire les delais de depistage,
de standardiser la pre-analyse et d'ameliorer la coordination entre les acteurs de soin.

\section{Methodologie du projet (Scrum)}
Pour cadrer le developpement de l'application, nous avons adopte une methode agile de type Scrum,
adaptee aux projets iteratifs en intelligence artificielle et en developpement logiciel.

Le travail a ete organise en sprints courts, avec des objectifs incrementaux :
\begin{itemize}
    \item preparation et exploration des donnees ;
    \item experimentation et amelioration des modeles ;
    \item integration backend/frontend/service IA ;
    \item validation fonctionnelle et tests de bout en bout.
\end{itemize}

Scrum a apporte trois benefices principaux dans ce projet :
\begin{itemize}
    \item \textbf{adaptabilite} : ajustement rapide des priorites selon les resultats experimentaux ;
    \item \textbf{visibilite} : suivi continu des taches et des livrables ;
    \item \textbf{qualite} : integration progressive et verification reguliere des fonctionnalites.
\end{itemize}

\section{Partie statistique}
\subsection{a) Dans le monde}
Les indicateurs epidemiologiques internationaux confirment l'ampleur du probleme :
\begin{itemize}
    \item plus de 347 millions de personnes vivent avec le diabete dans le monde ;
    \item la retinopathie diabetique concernerait plus de 93 millions de personnes ;
    \item environ 40\% a 45\% des personnes diabetiques presentent un stade de retinopathie.
\end{itemize}

Ces chiffres montrent que le depistage de la RD n'est pas uniquement un enjeu clinique,
mais aussi un enjeu de sante publique globale, en particulier dans les regions ou
l'acces aux specialistes reste limite.
\textit{Sources : WHO Fact Sheet Diabetes (2024), IDF Diabetes Atlas (2025).}

\subsection{b) En Tunisie}
En Tunisie, le diabete represente egalement un probleme majeur de sante publique,
avec une prevalence estimee autour de 16\% chez les adultes (20--79 ans) selon la serie
IDF/World Bank (2024). Cette progression s'accompagne d'un risque accru de complications
chroniques, dont la retinopathie diabetique.

Dans ce contexte, le depistage precoce reste un defi operationnel, notamment en raison :
\begin{itemize}
    \item de la concentration des specialistes dans les grands centres urbains ;
    \item des delais de consultation en ophtalmologie ;
    \item de l'augmentation continue du nombre de patients diabetiques.
\end{itemize}

L'integration d'une solution d'aide au depistage basee sur l'IA peut ainsi contribuer
a reduire les delais d'orientation et a ameliorer l'acces au diagnostic sur le territoire.
\textit{Source : World Bank Data -- Indicator SH.STA.DIAB.ZS (Tunisia, 2024).}

\section*{References web utilisees}
\addcontentsline{toc}{section}{References web utilisees}
\begin{itemize}
    \item World Health Organization (WHO). \textit{Diabetes -- Fact sheet}.\newline
    \url{https://www.who.int/news-room/fact-sheets/detail/diabetes}
    \item Centers for Disease Control and Prevention (CDC). \textit{Diabetes Basics}.\newline
    \url{https://www.cdc.gov/diabetes/about/index.html}
    \item International Diabetes Federation (IDF). \textit{IDF Diabetes Atlas}.\newline
    \url{https://diabetesatlas.org/}
    \item World Bank Data. \textit{Diabetes prevalence (\% of population ages 20 to 79) -- Tunisia}.\newline
    \url{https://data.worldbank.org/indicator/SH.STA.DIAB.ZS?locations=TN}
    \item NCBI Bookshelf (StatPearls). \textit{Diabetic Retinopathy}.\newline
    \url{https://www.ncbi.nlm.nih.gov/books/NBK560805/}
\end{itemize}

\section{Conclusion}
Ce chapitre a presente le contexte general du projet, en partant des bases medicales du diabete
et de la retinopathie diabetique, puis en decrivant les methodes diagnostiques actuelles et leurs limites.
L'analyse des contraintes terrain justifie l'interet d'une solution numerique intelligente.

La suite du rapport s'orientera vers l'etat de l'art des approches IA en imagerie retinienne,
les choix methodologiques retenus et la conception technique de la solution proposee.

% =========================

\chapter{Contexte general}
\label{chap:contexte_general}

\section{Introduction}
Le diabete sucré constitue aujourd'hui l'une des maladies chroniques les plus repandues dans le monde.
Parmi ses complications, la retinopathie diabetique (RD) est particulierement critique car elle peut evoluer
silencieusement pendant plusieurs annees avant l'apparition de symptomes visuels severes.
Sans depistage precoce, cette pathologie peut conduire a une perte de vision irreversible, voire a la cecite.

Sur le plan epidemiologique, la charge de morbidite est considerable. Les donnees de reference
internationales (CDC/OMS, reprises dans les ressources Kaggle de depistage de la RD) indiquent
que plus de 347 millions de personnes vivent avec le diabete dans le monde,
qu'environ 29.1 millions de personnes sont diabetiques aux Etats-Unis,
que la retinopathie diabetique concerne plus de 93 millions de personnes,
et qu'environ 40\% a 45\% des patients diabetiques presentent un certain stade de RD.
Ces chiffres confirment que la RD constitue un enjeu majeur de sante publique,
notamment dans les populations actives.
\textit{Sources : OMS (2024), CDC (2026), NCBI StatPearls (2023).}

Dans les structures de soins, la prise en charge repose encore majoritairement sur une lecture manuelle des images
de fond d'oeil par des specialistes. Cette approche est fiable, mais elle devient difficile a generaliser a grande
échelle en raison du volume de patients, des delais et du manque d'ophtalmologues dans certaines regions.

Dans ce cadre, l'intelligence artificielle (IA) apparait comme une solution d'assistance clinique pertinente :
elle permet d'automatiser une pre-analyse des images, de prioriser les cas urgents et de soutenir la decision
medicale sans remplacer l'expertise du medecin.

\section{Definition de la maladie}
\subsection{Diabete : rappel general}
Le diabete est une maladie metabolique caracterisee par une hyperglycemie chronique.
Cette hyperglycemie peut resulter d'un deficit de secretion d'insuline, d'une resistance a l'insuline,
ou d'une combinaison des deux. A long terme, l'exces de glucose sanguin provoque des atteintes vasculaires
et nerveuses qui touchent plusieurs organes (oeil, rein, coeur, systeme nerveux).

\subsection{Types de diabete}
\begin{itemize}
    \item \textbf{Diabete de type 1} : maladie auto-immune, destruction des cellules beta pancreatiques,
    dependance a l'insuline.
    \item \textbf{Diabete de type 2} : forme la plus frequente, associee a l'insulino-resistance et au mode de vie.
    \item \textbf{Diabete gestationnel} : hyperglycemie detectee pendant la grossesse.
    \item \textbf{Autres formes} : causes genetiques, endocriniennes ou medicamenteuses.
\end{itemize}

\subsection{Anatomie de l'oeil et zones d'interet}
Pour comprendre la retinopathie diabetique, il est essentiel de rappeler les structures oculaires impliquees
dans le diagnostic :
\begin{itemize}
    \item \textbf{Retine} : tissu neurosensoriel convertissant la lumiere en influx nerveux.
    \item \textbf{Macula} : zone centrale de la vision fine, critique pour la lecture et la precision visuelle.
    \item \textbf{Fovea} : centre de la macula, point de plus haute acuite visuelle.
    \item \textbf{Papille optique} : zone d'emergence du nerf optique.
    \item \textbf{Vaisseaux retiniens} : structures les plus affectees par les complications diabetiques.
\end{itemize}

\begin{figure}[H]
    \centering
    \includegraphics[width=0.72\textwidth]{images/ch1/anatomie_oeil.png}
    \caption{Anatomie simplifiee de l'oeil et principales zones d'interet pour le depistage de la RD.}
    \label{fig:anatomie_oeil_ch1}
\end{figure}

\subsection{Retinopathie diabetique : mecanisme et progression}
La retinopathie diabetique (RD) est une microangiopathie chronique de la retine associee
a un diabete mal controle ou ancien. L'hyperglycemie prolongee, associee a d'autres facteurs
de risque (hypertension arterielle, dyslipidemie, surcharge ponderale), fragilise l'endothelium
vasculaire et provoque des anomalies progressives de la microcirculation retinienne.

\subsubsection*{Lesions retiniennes elementaires}
Dans les images du fond d'oeil, la RD est principalement reconnue par la presence de lesions
vasculaires caracteristiques :
\begin{itemize}
    \item microanevrismes ;
    \item hemorragies intraretiniennes ;
    \item exsudats lipidiques ;
    \item exsudats cotonneux (zones d'ischemie nerveuse) ;
    \item anomalies veineuses (beading, irregularites) et ischemie ;
    \item neovascularisation dans les stades avances.
\end{itemize}

\subsubsection*{Stades cliniques : forme non proliferante et proliferante}
Sur le plan clinique, la progression suit deux grandes phases.

\paragraph{1) Retinopathie diabetique non proliferante (NPDR)}
Il s'agit du stade initial, le plus frequent au debut de l'evolution. Les petits vaisseaux
retiniens deviennent permeables, avec fuite de sang et de liquide. Cette fuite peut entrainer
un oedeme maculaire, cause majeure de baisse visuelle chez les personnes diabetiques.
Dans les stades non proliferants, la vision peut rester longtemps peu alteree, ce qui retarde
la consultation et renforce l'importance du depistage periodique.

\paragraph{2) Retinopathie diabetique proliferante (PDR)}
Lorsque l'ischemie retinienne devient importante (vaisseaux occlus), la retine declenche
une neovascularisation pathologique. Ces nouveaux vaisseaux sont fragiles, inefficaces
et susceptibles de saigner dans le vitre, provoquant flou visuel et corps flottants.
Le stade proliferant est associe a un risque eleve de complications severes.

\subsubsection*{Complications majeures des formes avancees}
\begin{itemize}
    \item \textbf{Hemorragie du vitre} : baisse visuelle brutale, parfois severe.
    \item \textbf{Fibrose retinienne} : formation de tissu cicatriciel pouvant deformer la retine.
    \item \textbf{Decollement tractionnel de retine} : urgence fonctionnelle menacant la vision centrale.
    \item \textbf{Glaucome neovasculaire} : obstruction de l'angle irido-corneen et elevation rapide
    de la pression intraoculaire.
\end{itemize}

\subsubsection*{Classification simplifiee exploitee dans le projet}
Pour la modelisation IA et l'evaluation, on retient une gradation en cinq niveaux :
\begin{enumerate}
    \item \textbf{Absence de RD} ;
    \item \textbf{RD non proliferante legere} ;
    \item \textbf{RD non proliferante moderee} ;
    \item \textbf{RD non proliferante severe} ;
    \item \textbf{RD proliferante} (risque eleve de perte visuelle).
\end{enumerate}

\begin{figure}[H]
    \centering
    \includegraphics[width=0.85\textwidth]{images/ch1/stades_retinopathie.png}
    \caption{Exemple visuel des stades de progression de la retinopathie diabetique.}
    \label{fig:stades_rd_ch1}
\end{figure}

\paragraph{Lien avec le depistage automatise}
Les informations precedentes montrent qu'une detection rapide des lesions est essentielle
pour eviter la progression vers les formes proliferantes. C'est exactement l'objectif
des systemes d'aide au depistage automatise : analyser les images de fond d'oeil,
estimer un grade de severite et orienter en priorite les patients a risque eleve.

\section{Methodes de diagnostic existantes}
\subsection{Examen clinique du fond d'oeil}
L'examen du fond d'oeil est la methode standard pour observer les lesions retiniennes.
Il peut etre realise avec ou sans dilatation pupillaire, selon le contexte clinique.

\subsection{Retinographie couleur}
La retinographie permet de capturer des images numeriques du fond d'oeil.
Elle facilite l'archivage, la tele-expertise et le suivi longitudinal des patients.

\subsection{Tomographie en coherence optique (OCT)}
L'OCT fournit des coupes de la retine a haute resolution.
Cette technique est particulierement utile pour detecter l'oedeme maculaire diabetique.

\subsection{Angiographie retinienne}
L'angiographie a la fluoresceine visualise la perfusion vasculaire et les zones de fuite.
Bien qu'informative, elle est plus invasive et plus couteuse.

\subsection{Tele-ophtalmologie}
La tele-ophtalmologie permet de capturer localement les images puis de les transmettre
a un expert a distance. Cette approche augmente l'accessibilite au depistage.

\begin{table}[H]
\centering
\caption{Comparaison des principales methodes de diagnostic de la retinopathie diabetique.}
\label{tab:comparaison_diagnostic_ch1}
\begin{tabular}{|p{3.8cm}|p{3.2cm}|p{3.6cm}|p{3.6cm}|}
\hline
\textbf{Methode} & \textbf{Atouts} & \textbf{Limites} & \textbf{Contexte d'usage} \\
\hline
Fond d'oeil clinique & Methode de reference, pratique courante & Dependance a l'expertise humaine & Depistage et suivi \\
\hline
Retinographie & Trace numerique, partage a distance & Qualite variable selon acquisition & Campagnes de depistage \\
\hline
OCT & Tres haute precision structurelle & Cout eleve, equipement specialise & Oedeme maculaire, cas complexes \\
\hline
Angiographie & Analyse vasculaire detaillee & Invasive, moins accessible & Exploration avancee \\
\hline
Tele-ophtalmologie & Couvre les zones eloignees & Delai d'interpretation possible & Reseaux de soins distribues \\
\hline
\end{tabular}
\end{table}

\section{Prediction du diabete par l'intelligence artificielle}
Dans le contexte actuel, le depistage manuel de la retinopathie diabetique reste efficace
mais il est limite par la charge de travail des specialistes, la variabilite de lecture,
les delais de restitution et les inegalites d'acces aux soins.
L'intelligence artificielle repond a cette problematique en fournissant une aide rapide
et standardisee pour analyser les images du fond d'oeil, estimer un grade de severite,
et prioriser les cas urgents.

L'objectif n'est pas de remplacer l'ophtalmologue, mais de l'assister dans la decision,
de reduire le retard diagnostique et d'ameliorer la qualite globale du depistage.

\subsubsection*{1.4.7 Solution avec IA et concept de l'application}
La solution proposee dans ce projet repose sur une application d'aide au depistage
qui combine une interface clinique (centre/medecin), un backend de gestion des examens
et un service IA d'analyse d'images retiniennes.

Le concept de l'application est le suivant :
\begin{enumerate}
    \item acquisition de l'image du fond d'oeil par le centre ;
    \item envoi securise de l'image vers le serveur d'analyse ;
    \item prediction automatique du grade de retinopathie ;
    \item generation d'une explication visuelle (Grad-CAM) ;
    \item notification du medecin, avec priorisation des cas severes.
\end{enumerate}

Cette architecture permet de reduire les delais de depistage,
de standardiser la pre-analyse et d'ameliorer la coordination entre les acteurs de soin.

\section{Methodologie du projet (Scrum)}
Pour cadrer le developpement de l'application, nous avons adopte une methode agile de type Scrum,
adaptee aux projets iteratifs en intelligence artificielle et en developpement logiciel.

Le travail a ete organise en sprints courts, avec des objectifs incrementaux :
\begin{itemize}
    \item preparation et exploration des donnees ;
    \item experimentation et amelioration des modeles ;
    \item integration backend/frontend/service IA ;
    \item validation fonctionnelle et tests de bout en bout.
\end{itemize}

Scrum a apporte trois benefices principaux dans ce projet :
\begin{itemize}
    \item \textbf{adaptabilite} : ajustement rapide des priorites selon les resultats experimentaux ;
    \item \textbf{visibilite} : suivi continu des taches et des livrables ;
    \item \textbf{qualite} : integration progressive et verification reguliere des fonctionnalites.
\end{itemize}

\section{Partie statistique}
\subsection{a) Dans le monde}
Les indicateurs epidemiologiques internationaux confirment l'ampleur du probleme :
\begin{itemize}
    \item plus de 347 millions de personnes vivent avec le diabete dans le monde ;
    \item la retinopathie diabetique concernerait plus de 93 millions de personnes ;
    \item environ 40\% a 45\% des personnes diabetiques presentent un stade de retinopathie.
\end{itemize}

Ces chiffres montrent que le depistage de la RD n'est pas uniquement un enjeu clinique,
mais aussi un enjeu de sante publique globale, en particulier dans les regions ou
l'acces aux specialistes reste limite.
\textit{Sources : WHO Fact Sheet Diabetes (2024), IDF Diabetes Atlas (2025).}

\subsection{b) En Tunisie}
En Tunisie, le diabete represente egalement un probleme majeur de sante publique,
avec une prevalence estimee autour de 16\% chez les adultes (20--79 ans) selon la serie
IDF/World Bank (2024). Cette progression s'accompagne d'un risque accru de complications
chroniques, dont la retinopathie diabetique.

Dans ce contexte, le depistage precoce reste un defi operationnel, notamment en raison :
\begin{itemize}
    \item de la concentration des specialistes dans les grands centres urbains ;
    \item des delais de consultation en ophtalmologie ;
    \item de l'augmentation continue du nombre de patients diabetiques.
\end{itemize}

L'integration d'une solution d'aide au depistage basee sur l'IA peut ainsi contribuer
a reduire les delais d'orientation et a ameliorer l'acces au diagnostic sur le territoire.
\textit{Source : World Bank Data -- Indicator SH.STA.DIAB.ZS (Tunisia, 2024).}

\section*{References web utilisees}
\addcontentsline{toc}{section}{References web utilisees}
\begin{itemize}
    \item World Health Organization (WHO). \textit{Diabetes -- Fact sheet}.\newline
    \url{https://www.who.int/news-room/fact-sheets/detail/diabetes}
    \item Centers for Disease Control and Prevention (CDC). \textit{Diabetes Basics}.\newline
    \url{https://www.cdc.gov/diabetes/about/index.html}
    \item International Diabetes Federation (IDF). \textit{IDF Diabetes Atlas}.\newline
    \url{https://diabetesatlas.org/}
    \item World Bank Data. \textit{Diabetes prevalence (\% of population ages 20 to 79) -- Tunisia}.\newline
    \url{https://data.worldbank.org/indicator/SH.STA.DIAB.ZS?locations=TN}
    \item NCBI Bookshelf (StatPearls). \textit{Diabetic Retinopathy}.\newline
    \url{https://www.ncbi.nlm.nih.gov/books/NBK560805/}
\end{itemize}

\section{Conclusion}
Ce chapitre a presente le contexte general du projet, en partant des bases medicales du diabete
et de la retinopathie diabetique, puis en decrivant les methodes diagnostiques actuelles et leurs limites.
L'analyse des contraintes terrain justifie l'interet d'une solution numerique intelligente.

La suite du rapport s'orientera vers l'etat de l'art des approches IA en imagerie retinienne,
les choix methodologiques retenus et la conception technique de la solution proposee.
